\documentclass[a4paper,12pt]{extarticle}
\usepackage[table,xcdraw]{xcolor}
\usepackage{pgfplots}
\usepackage[unicode, pdftex]{hyperref}
\usepackage[T2A]{fontenc}
\usepackage[utf8]{inputenc}
\usepackage[english,russian]{babel}
\usepackage{amsmath,amsfonts,amsthm, mathtools}
\usepackage{amssymb}
\usepackage{dsfont}
\usepackage{graphicx}
\graphicspath{{images/}}
\usepackage{wrapfig}
\usepackage{floatrow}
\usepackage{setspace}
\usepackage[headings]{fancyhdr}
\usepackage[margin=10pt,font={small,stretch=0,9},labelfont=bf,labelsep=period,
justification=centerlast]{caption}
\usepackage{array,tabularx,tabulary,booktabs}
\RequirePackage{multirow}
\RequirePackage{multicol}
\RequirePackage{longtable}
\usepackage{indentfirst}
\usepackage{hyperref}
\usepackage{icomma}
\newcommand*{\hm}[1]{#1\nobreak\discretionary{}
	{\hbox{$\mathsurround=0pt #1$}}{}}
\usepackage{soulutf8} 
\usepackage{etoolbox} 
\usepackage{geometry}
\geometry{top=20mm}
\geometry{bottom=20mm}
\geometry{left=20mm}
\geometry{right=20mm}
\DeclareMathOperator{\sgn}{\mathop{sgn}}
\usepackage[OT1]{fontenc}
\usepackage{amsmath}
\usepackage{amsfonts}
\usepackage{amssymb}
\usepackage{wasysym}
\usepackage{wrapfig}
\usepackage{physics}
\usepackage[normalem]{ulem}
\graphicspath{{pictures/}}
\DeclareGraphicsExtensions{.png,.jpg,.jpeg}
\usepackage{indentfirst}
\usepackage[T2A]{fontenc}
\usepackage{subfigure}
\usepackage{fancyhdr}
\usepackage{caption}

\usepackage{tabularx}
\usepackage{floatrow}
\usepackage{multicol}
\usepackage{lastpage}

\pgfplotsset{
  MyStyle1/.style={
    axis lines = middle,
    enlargelimits = true,
    grid=both,
    grid style={line width=.1pt, draw=gray!10},
    major grid style={line width=.2pt,draw=gray!50},
    minor tick num=5,
    enlargelimits=0.05,
    ticklabel style={font=\tiny,fill=white},
    xlabel style={at={(ticklabel* cs:1)},anchor=north west},
    ylabel style={at={(ticklabel* cs:1)},anchor=south east},
  }
}
\def\Slope{1.5}
\def\Intercept{-20}


% Настройка отчёркиваний и нумерации
\usepackage{fancyhdr}
\pagestyle{fancy}
\fancyhf{}

\fancyhead[L]{\nouppercase{\leftmark\hfill}}
\fancyfoot[RO, RE]{Страница \thepage \ из \pageref{LastPage}}
\renewcommand{\headrulewidth}{0.5pt}
\renewcommand{\footrulewidth}{0.5pt}

\usepackage{titlesec}
\titlelabel{\thetitle.\quad}

\pgfplotsset{compat=1.18}
\begin{document}

\thispagestyle{empty}
\begin{center}
	\textit{Федеральное государственное автономное образовательное\\ учреждение высшего образования }
	
	\vspace{0.5ex}
	
	\textbf{«Московский физико-технический институт\\ (национальный исследовательский университет)»}
\end{center}

\vspace{10ex}
\begin{figure}[!h]
  \begin{center}
    
\includegraphics[width = 0.8\textwidth]{MIPT.png}
  \end{center}
\end{figure}

\begin{center}
	
	\textbf{Лабораторная работа № 4.7.3}
	
	\vspace{1ex}
	
	по курсу общей физики
	
	на тему:
	
	\textbf{\textit{<<Поляризация>>}}
	
	\vspace{20ex}
	
	\begin{flushright}
		\noindent
		\textit{Работу выполнил:}\\  
		\textit{Легких Никита \\(группа Б05-301)}
	\end{flushright}
	\vfill
	Долгопрудный \\ \today
	\setcounter{page}{1}
\end{center}
\newpage

% Конец шаблона
\newpage

\textbf{Цель работы:} ознакомление с методами получения и анализа поляризованного света.

\textbf{В работе используются:} оптическая скамья с осветителем; зелёный светофильтр; два поляроида (цена деления держателя 1$^{\circ}$); чёрное зеркало; полированная эбонитовая пластинка; стопа стеклянных пластинок; слюдяные пластинки разной толщины; пластинки в 1/4 и 1/2 длины волны; пластинка в одну длину волны для зелёного света (пластинка чувствительного
оттенка).

\section*{Ход работы}
\subsection*{Определение разрешенных направлений поляроидов}
\begin{wrapfigure}{r}{0.4\textwidth}
    \includegraphics[width=1\textwidth]{pics/1.png}
    \caption{}
\end{wrapfigure}
С помощью черного зеркала определим разрешенное направление первого поляроида. Для этого установим осветитель $S$, поляроид $P_1$ и черное зеркало в соответствии со схемой на рис.\ref{schema_1} и, поворачивая поляроид вокруг направления луча, а черное зеркало вокруг вертикальной оси, добьемся наименьшей яркости. С помощью лимба поляроида получим, что разрешенное направление поляроида соответствует $(312 \pm 1)^{\circ}$ на лимбе.\\
Заменив черное зеркало на рис.1 на второй поляроид $P_2$ и добившись минимальной яркости луча, получим, что разрешенное направление второго поляроида соответствует $(5 \pm 1)^{\circ}$ на лимбе.
\subsection*{Определение угла Брюстера для эбонита}
Заменим черное зеркало на рис.1 на эбонитовую пластину с круговой шкалой и повернем эбонитовое зеркало вокруг вертикальной оси так, чтобы его плоскость была перпендикулярна лучу. Получим, что начало отсчета соответствует $(353 \pm 1)^{\circ}$ на круговой шкале.\\
Теперь найдем угол поворота эбонита $\varphi_{B}$, при котором интенсивность отраженного луча минимальна. Он соответствует $(50 \pm 1)^{\circ}$ на круговой шкале. Следовательно, сам угол Брюстера равен: $\varphi_{B} = (57 \pm 2)^{\circ}$.\\
Исходя из полученного значения, рассчитаем показатель преломления эбонита:\\
\begin{equation*}
    n = \tan \varphi_{B} = (1,54 \pm 0,12) 
\end{equation*}
Погрешность считаем по формуле: $\Delta n = \frac{\Delta \varphi}{\sin^2 \varphi}$
Заметим, что при повторе измерений с добавлением зеленого светофильтра $Ф$ мы получаем аналогичные результаты.
\subsection{Исследование стопы}
\begin{wrapfigure}{r}{0.4\textwidth}
    \includegraphics[width=1\textwidth]{pics/2.png}
    \caption{}
\end{wrapfigure}
Поставим стопу стеклянных пластинок вместо эбонитового зеркала и подберем для нее такое положение, при котором свет падает на стопу под углом Брюстера. Он соответствует $(40 \pm 1)^{\circ}$. Далее осветим стопу неполяризованным светом и, рассматривая через поляроиды свет (рис.2, отраженный от стопы, определим ориентацию вектора $\mathbf{E}$ в отраженном луче, затем повторим это для преломленного луча. Получим, что плоскости поляризации у отраженного и преломленного лучей взаимно перпендикулярны: угол на лимбе $P_1 = (40 \pm 1)^{\circ}$, $P_2 = (275 \pm 1)^{\circ}$. Преломленные лучи горизонтальные, отраженные – вертикальные. 
\subsection*{Определение главных плоскостей двоякопреломляющих пластин}
\begin{wrapfigure}{r}{0.4\textwidth}
    \includegraphics[width=1\textwidth]{pics/3.png}
    \caption{}
\end{wrapfigure}
Поставим кристаллическую пластинку между скрещенными поляроидами $P_1$ и $P_2$ и, вращая пластинку вокруг направления луча и наблюдая за интенсивностью света, проходящего сквозь второй поляроид, определим, при каком условии главные направления пластинки совпадают с разрешенными направлениями поляроидов. Повторим опыт для второй пластинки. Минимумы и максимумы интенсивности чередуются через 45°, главные плоскости пластин совпадают с разрешенными направлениями поляроидов при максимальной интенсивности. Минимумы: для первой $(10 \pm 1)^{\circ}$, для второй $(322 \pm 1)^{\circ}$.
\subsection*{Выделение пластин $\lambda/2$ и $\lambda/4$}
Добавим к схеме, изображенной на рис.3, зеленый фильтр, установим разрешенное направление поляроида горизонтально, а главные направления исследуемой пластинки — под углом 45° к горизонтали. С помощью второго поляроида установим, какую поляризацию имеет свет, прошедший пластинку.\\
Для первой получим, что при вращении интенсивность света не меняется, а также он остается зеленым, следовательно, поляризация круговая и пластинка соответствует $\lambda/4$. Для второй — что интенсивность света меняется при вращении, а также он приобретает красно-фиолетовые оттенки, следовательно, поляризация линейная и пластинка соответствует $\lambda/2$.
\subsection*{Определение направлений большей и меньшей скорости в пластинке $\lambda/4$}
\begin{wrapfigure}{r}{0.4\textwidth}
    \includegraphics[width=1\textwidth]{pics/4.png}
    \caption{}
\end{wrapfigure}
Поставим между скрещенными поляроидами пластинку чувствительного оттенка ($\lambda$ для зеленого света), имеющую вид стрелки. Световой вектор, ориентированный вдоль направления стрелки, проходит с большей скоростью, перпендикулярный — с меньшей. Установим разрешенное направление первого поляроида горизонтально и убедимся с помощью второго поляроида, что эта пластинка
не меняет поляризацию зеленого света в условиях предыдущего опыта.\\
Уберем зеленый фильтр и поставим между скрещенными поляроидами пластинку $\lambda$ (стрелка под углом 45° к разрешенным направлениям поляроидов). Глядя сквозь второй поляроид на стрелку, убедимся, что она имеет пурпурный цвет (зеленый свет задерживается вторым поляроидом, а красная и синяя компоненты проходят). Добавим к схеме пластинку $\lambda/4$ (рис.4), главные направления которой совпадают с главными направлениями пластины $\lambda$ и ориентированы под углом 45° к разрешенным направлениям скрещенных поляроидов.\\
При повороте рейтера со стрелкой на 180° вокруг вертикальной оси цвет стрелки меняется от зелено-голубого до оранжево-желтого. Этот эффект объясняется следующим образом: при совпадении быстрой оси стрелки и главного направления пластинки суммарная разность хода для зеленого света составит $5\lambda/4$, что приблизительно равно длине волны в красном спектре, поэтому красный свет будет задерживаться поляроидами и стрелка окрашивается в зелено-голубой цвет, при совпадении же медленной оси с главным направлением пластинки разность хода $3\lambda /4$ будет примерно равна длине волны для фиолетово-голубой части спектра, поэтому в результате мы видим оранжево-желтый цвет стрелки.\\
Отметим, что положение быстрой оси пластинки соответствует $(205 \pm 1)^{\circ}$.
\subsection*{Интерференция поляризованных лучей}
Исследуем интерференцию поляризованных лучей. Для этого расположим между скрещенными поляроидами мозаичную слюдяную пластинку. Она собрана из 4-х узких полосок слюды, лежащих по сторонам квадрата (две полоски «толщиной» $ \lambda / 4 $ и по одной — $ \lambda / 2 $ и $ 3 \lambda / 4 $). В центральном квадратике слюды нет. Главные направления всех пластинок ориентированы параллельно сторонам квадрата.\\
Вращая пластинку, пронаблюдаем за изменениями в отдельном квадратике. У нас изменяется интенсивность.\\
Не трогая пластинки, начнем вращать второй поляроид. Отличие в том, что теперь изменяется цвет.
\subsection*{Определение направления вращения светового вектора в эллиптически поляризованной волне}
\begin{figure}[H]
	\centering
	\includegraphics[width=0.7\textwidth]{pics/5.png}
    \caption{Эллипс поляризации и графики синусоид}
\end{figure}
Снова поставим зеленый фильтр, а за ним между скрещенными поляроидами -- пластинку произвольной толщины $( \lambda / 4 )$.\\
Получим эллиптически-поляризованный свет. Для этого установим разрешённое направление первого поляроида под углом 10°--20° к горизонтали так, чтобы вектор \textbf{E} падающего на пластинку света был расположен в первом квадранте. Установим разрешённое направление второго поляроида вертикально и, вращая пластинку, найдем минимальную интенсивность света, прошедшего второй поляроид. Вращая второй поляроид, убедитесь, что свет поляризован эллиптически, а не линейно. Таким образом, получим эллипс поляризации с вертикально ориентированной малой осью.\\
Для определения направления вращения светового вектора в эллипсе установим между поляроидами дополнительную пластинку $ \lambda / 4 $ с известными направлениями «быстрой» и «медленной» осей, ориентированными по осям эллипса поляризации анализируемого света. В этом случае вектор \textbf{E} на выходе будет таким, как если бы свет прошел две пластинки $ \lambda / 4 $: свет на выходе из второй пластинки будет линейно поляризован. Если пластинки поодиночке дают эллипсы, вращающиеся в разные стороны, то поставленные друг за другом, они скомпенсируют разность фаз, и вектор \textbf{E} на выходе останется в первом и третьем квадрантах. Если же световой вектор перешёл в смежные квадранты, значит, эллипсы вращаются в одну сторону.\\
После второго поляроида интенсивность света максимальна в первом квадранте. Значит, две пластины компенсируют друг друга, световой вектор остался в том же квадранте, эллипсы вращаются в разные стороны.
\section*{Выводы}
В ходе лабораторной работы мы ознакомились с методами получения и анализа поляризованного света и получили следующие результаты:
\begin{itemize}
\item Определены разрешённые направления поряроидов --- для первого поляроида разрешённое направление горизонтальное, на лимбе $(312 \pm 1)^{\circ}$, для второго поляроида --- вертикальная волна, на лимбе $(5 \pm 1)^{\circ}$;
\item Был определен показатель преломления эбонита $n = (1,54 \pm 0,12)$, что сходится с табличным значением в рамках погрешности $n_{\text{табл}} = 1,6-1,7$;
\item Получили, что в стопе стеклянных пластинок при угле Брюстера преломленные лучи горизонтальные, а отраженные --- вертикальные;
\item Для двоякопреломляющих пластин определены главные направления --- минимумы и
максимумы интенсивности чередуются через $45^{\circ}$, главные плоскости пластин совпадают с разрешенными направлениями поляроидов при максимальной интенсивности. Минимумы: для первой $(10 \pm 1)^{\circ}$, для второй $(322 \pm 1)^{\circ}$;
\item Выяснили, что из двух имеющихся пластинок первая из них --- пластинка $\lambda/4$, а вторая --- $\lambda/2$;
\item «Быстрая» ось пластинки $\lambda/4$ совпадает с «быстрой» осью пластинки $\lambda$ и ориентирована под углом $45^{\circ}$ к разрешенным направлениям скрещенных поляроидов. При повороте рейтера со стрелкой на $180^{\circ}$ вокруг вертикальной оси получаем медленную ось.
\item При изучении интерференции поляризованных лучей было получено, что при вращении пластинки в отдельном квадратике изменяется интенсивность, а при вращении второго поляроида --- изменяется цвет;
\item Было определено направление вращения светового вектора в эллиптически поляризованной волне после пластинки $\lambda/4$: оно компенсирует разность фаз со второй пластинкой, «быстрая» и «медленная» оси которой известны. Таким образом, эллипс первой пластинки вращается в противоположную сторону.
\end{itemize}

\end{document}